\chapter{The analytic encoder}\label{ch:encoder}
\chapterbrief
{a Schr\"odinger-type equation with a variable potential $q(x)$ \emph{known}, written $P(\partial_t,L_q)u=0$ with $L_q=-\partial_x^2+q$ and $L_0=-\partial_x^2$, on a symmetric interval $J$.}
{the variable potential is a nuisance: the theory of Phases~1--5 applies to $L_0$. Can one change the unknown by an explicit, invertible, stable map so that the potential disappears, without machine learning? Such a map is an \emph{encoder}.}
{supplies Step~5 of the extended algorithm (Chapter~\ref{ch:extalgo}): the encoder is a \emph{translation} to constant coefficients, not a source of information. The proofs are Dim.~A44--A47 and B79.}

\section{What an encoder is}

A \emph{structured encoder} is a linear map $W$ acting in $x$ only, with (E1) finite-dimensional kernel on its own domain, (E2$'$) $WL_q=(L_0+c)W$ for a constant $c$, and (E3) explicit. The conditions are all needed: without (E1) the zero map qualifies, and without (E3) the interaction-picture map $U_0(t)U_V(t)^{-1}$ does.

\begin{proposition}{A constant potential has no strict differential intertwiner \led{ENC-K0} \proved}\label{prop:enck0}
For $q\equiv c\ne0$ no differential operator of any order satisfies $WL_c=L_0W$. With the relaxed requirement $WL_c=(L_0+c)W$ the identity works.
\dimref{ENC-K0}
\end{proposition}

\section{The Volterra transformation operator}

Let $K$ solve the Goursat problem $K_{xx}-K_{ss}=q(x)K$ on the triangle $|s|\le|x|\le a$, with $K(x,x)=\tfrac12\int_0^xq$ and $K(x,-x)=0$, and put
\[
Tf(x)=f(x)+\int_{-x}^xK(x,s)f(s)\dd s,\qquad W=T^{-1}.
\]

\begin{theorem}{The symmetric Volterra encoder \led{ENC-1, ENC-U1} \proved}\label{thm:enc1}
For $q\in C^1$ the problem has a unique solution, and for every $f\in C^2(J)$, without any boundary condition, $L_qTf=TL_0f$ with $(Tf)(0)=f(0)$, $(Tf)'(0)=f'(0)$. $T$ is bijective on $C$, $C^2$ and (for $q\in C^\infty$) $C^\infty$, with $\|K\|\le\frac{aQ}2e^{a\sqrt Q}$ and a Lipschitz dependence on $q$. Among all Volterra intertwiners the one preserving Cauchy data at the base point is unique; it depends on the base point already for constant $q$.
\dimref{ENC-1}
\end{theorem}

\begin{intuition}
In characteristic coordinates the kernel solves a Volterra integral equation by successive approximations, with the bound $(Q|\xi\eta|/4)^n/(n!)^2$. Volterra operators with a triangular kernel have no spectrum, so $T$ is invertible, and the Goursat conditions are exactly the ones that make every boundary term in the Leibniz expansion of $L_qTf-TL_0f$ cancel.
\end{intuition}

\begin{corollary}{One encoder for every equation \led{ENC-C1} \proved}\label{cor:encc1}
For $q\in C^\infty$ and every polynomial $P$: $P(\partial_t,L_q)u=0\iff P(\partial_t,L_0)(Wu)=0$. The same $W$ serves Schr\"odinger, heat, wave and Klein--Gordon.
\dimref{ENC-C1}
\end{corollary}

\begin{example}{Constant potential}
For $q\equiv c$ the transformation maps the exponential $e^{ik\cdot}$ to $e^{ikx_0}[\cos\mu(x-x_0)+ik\sin\mu(x-x_0)/\mu]$ with $\mu^2=k^2-c$: the eigenfunction of $L_c$ with the Cauchy data of $e^{ikx}$ at $x_0$.
\end{example}

\section{What no encoder can do}

\begin{theorem}{No time-independent encoder for a time-dependent potential \led{ENC-P0} \pw}\label{thm:encp0}
Let $W$ be linear with finite kernel and $W(-\partial_x^2+V(t,\cdot))f=(-\partial_x^2+c(t))Wf$ for all $t$. Then $V(t,x)-V(s,x)=c(t)-c(s)$: the potential must be autonomous up to a function of time. The reason is that otherwise infinitely many independent functions $m\varphi_k$ would lie in $\ker W$. Encoders for a genuinely time-dependent potential exist but depend on time (interaction pictures), and a pointwise representative with the identity time change exists exactly for $V=\beta(t)x+\gamma(t)$.
\dimref{ENC-P0b}
\end{theorem}

\begin{theorem}{Blind encoders and identifiability \led{ENC-N1$'$} \proved}\label{lem:encn1}
A blind encoder, one whose output does not depend on which potential produced the data, exists iff the coherence condition (C) holds for the chosen family; it is implied by identifiability and is strictly weaker (a witness: the data model that observes only a phase). The correct statement is: blind encoder $\iff$ identifiability modulo the gauge group of the data model.
\dimref{ENC-N1p}
\end{theorem}

\begin{caution}[title=An earlier claim that turned out false]
The synthesis of the first version, ``the exceptions change the time'', was false as written: Avron--Herbst and the whole class $V=\beta(t)x+\gamma(t)$ have a pointwise representative with the identity time change. Likewise ``a blind encoder exists iff $W_V^{-1}v$ is injective'' is false (ENC-N1).
\end{caution}

\section{Other results of this phase}

\begin{itemize}
\item \textbf{Darboux encoders}: a differential $W$ with $WL_q=(L_0+c)W$ is a Darboux chain on each component of a complement of a countable closed set; for rational $V\to0$ it exists iff $V$ is an iterated Darboux transform of $0$ (Adler--Moser) \led{ENC-K1$'$, ENC-K1}. They are unbounded with a non-trivial kernel, so they lose modes.
\item \textbf{The interaction term with a potential}: $\|E^{\rm mult}f\|\le\frac{t^2}2(\|V''\|\|f\|+2\|V'\|\|f'\|)+\frac{2t^3}3\|V'\|^2\|f\|$ with the optimal constant $2/3$, and $E^{\rm mult}\equiv0$ iff $V$ is constant \led{ENC-A9}.
\item \textbf{Faddeev decomposition} on an interval and \textbf{Foldy--Wouthuysen} for the free Dirac $1+1$ are decompositions, not encoders \led{FAD-1, FW-1}.
\end{itemize}

\begin{takeaway}
Removing a variable potential is possible by an explicit stable Volterra transformation when the potential is known, and the same map works for every equation in $(\partial_t,L_q)$. When the potential is unknown the encoder gives no information, and for a time-dependent potential no time-independent encoder exists.
\end{takeaway}
